% Packets
\title{Author}
\author{Author}
\date{13/06/2023}

% Fonts
\usepackage[T1]{fontenc}
\usepackage{lmodern}	%	\normalfont	% Fuentes Latin modern
\usepackage[utf8]{inputenc}	% Distribución que permite formato de texto con tíldes, etc.
\usepackage[spanish]{babel}	% Español
\usepackage{booktabs}

% Chapters, Heading, footnotes and captions
\usepackage{fancyhdr}			% Paquete para notas en la Header y Foot de página
\usepackage{emptypage}			% Quita los encabezados en las páginas en blanco
\usepackage[pdftex,
            pdfauthor={Author},
            pdftitle={Report TP},
            pdfsubject={Reporte},
            %pdfkeywords={Cryogenic, photodetector, HgCdTe - MCT, dark current, cooling capacity, cryocooler},
            pdfproducer={Documentclass: "Book" with "Latex" in "Ubuntu-Linux"},
            pdfcreator={Commands: "PDFLaTeX" and "BibTex"},
	colorlinks=true, citecolor=black, linkcolor=black, urlcolor=black]{hyperref}
\usepackage[margin=5pt,font=footnotesize,labelfont=bf,labelsep=endash,textfont=it]{caption}	% Manage caption size
\usepackage{fnpos}				% Para usar \footnote{•}
\usepackage{ftnxtra}			% Para usar \footnote{•} dentro de un \caption{•}
\usepackage{lscape}				% Entorno para rotar contenido en la página
\usepackage[none]{hyphenat}		% Para que no se corte ninguna palabra y la continúe en el renglón de abajo.

% Equations, figs and tables
\usepackage{float}			% Reconoce campos de tablas, figuras, etc.
\usepackage{graphicx}		% Gráficos
\usepackage{subfig}			% Subfiguras
	\captionsetup[subfigure]{subrefformat=simple,labelformat=simple,listofformat=subsimple} % Conf. letra asoaciada a la sub-figura debajo de cada una.
	\renewcommand\thesubfigure{(\alph{subfigure})} % Conf. \ref de subfloats.
\usepackage{wrapfig}			 % Para colocar imágenes que sean rodeadas por texto
\usepackage{amssymb}			 % Símbolos especiales
	\usepackage{pifont}		 % Más símbolos que se agregan como \ding{##}
\usepackage{bm}				     % Fuente matemática: bolt + italic
\usepackage{amsmath,amsthm,wasysym}     % Paquete matemático - símbolos
	\usepackage{amsfonts}	 % Fuentes matemáticas
\usepackage{multirow}	% Unir celdas fila en las tablas
\usepackage{array}		% Permite fijar ancho de columnas en tablas
% La siguiente configuración permite fijar anchos de columnas alineadas directamente:
	\newcolumntype{L}[1]{>{\raggedright\let\newline\\\arraybackslash\hspace{0pt}}m{#1}}
	\newcolumntype{C}[1]{>{\centering\let\newline\\\arraybackslash\hspace{0pt}}m{#1}}
	\newcolumntype{R}[1]{>{\raggedleft\let\newline\\\arraybackslash\hspace{0pt}}m{#1}}
\usepackage{paralist}	% Permite usar \begin{compactitem} para un entorno itemize con espacios reducidos

% Anexxes:
\usepackage[toc,page]{appendix}					% Paquete para Apéndices (Anexos)
	\renewcommand{\spanishappendixname}	{Anexo} 	% Reemplaza la palabra Apéndice en cada uno
	\renewcommand{\appendixtocname}		{Anexos}	% Nombre del encabezado de Anexos en TOC
	\renewcommand{\appendixpagename}		{ANEXOS}	% Nombre de la portada de "ANEXOS"

% Bibliografía:
\usepackage[nottoc,notlot,notlof]{tocbibind}	% Para incluir la bibliografía en el índice

% Definición de colores a usar en la configuración de otros paquetes:
\usepackage{color}
\definecolor{gray97}{gray}{.97}
\definecolor{gray75}{gray}{.75}
\definecolor{gray45}{gray}{.45}
\setcounter{MaxMatrixCols}{20}	% 20 columnas. (Aceptar matrices con más de 10 columnas)

% Para tener referencia a la última página:
\usepackage[export]{adjustbox}
\usepackage{lastpage}
\usepackage{enumerate}	% Paquete con estilos de enumeración

% Dimensiones:
\usepackage{geometry}
\geometry{
	a4paper,
	margin=1in
}

% Interlineado:
\usepackage{setspace}
% Opciones: \onehalfspace; \doublespacing; \singlespace; \spacing{Nº}
\onehalfspace
\raggedbottom

% Enumeración de ítems:
\usepackage{blindtext}			% Habilita el uso de \begin{enumerate}
\usepackage[inline]{enumitem}	% Para poder enumerar dentro del texto con distintos formatos

% Estilo de capítulos:
\usepackage{titlesec} % (la versión del 2016/03/15 no numera "\section{•}")
%**********************************
%%%%%% CORRECCIÓN DEL BUG: %%%%%%
%**********************************
	\usepackage{etoolbox}
		\makeatletter
		\patchcmd{\ttlh@hang}{\parindent\z@}{\parindent\z@\leavevmode}{}{}
		\patchcmd{\ttlh@hang}{\noindent}{}{}{}
		\makeatother
%**********************************
% Configuración de capítulos:
\titleformat{\chapter}[display]
  {\bfseries\Huge}
  {\vspace{-1.5cm}\filright\MakeUppercase{\Huge\chaptertitlename} \fontsize{40}{40}\selectfont\thechapter}
  {1ex}
  {\titlerule\vspace{1ex}\filleft}
  [\vspace{1ex}\titlerule]
  \titlespacing*{\chapter}{0pt}{0pt}{20pt} % Espacio antes del título del "capítulo".

%%%%%%%%%%%%%%%%%%%%%%%%%%%%%%%%%%%%%
%%%% ENCABEZADO Y PIÉ DE PÁGINAS %%%%%
\newcommand{\EncYpiepaGpreambulo}{
% Estilo:
%%	Definición de la página inicial de cada capítulo:
% Configuración de la línea divisoria:
\fancypagestyle{plain}{
\renewcommand{\headrulewidth}{0.0 pt}
\renewcommand{\footrulewidth}{0.0 pt}
% Encabezado superior:
\fancyhead[L]{}
\fancyhead[C]{}
\fancyhead[R]{}
% Encabezado inferior:
\fancyfoot[L]{}
\fancyfoot[C]{\thepage}
\fancyfoot[R]{}
}
%%	Definición de las siguientes páginas:
\pagestyle{fancy}
% Configuración de la línea divisoria:
	\renewcommand{\headrulewidth}{0.0 pt}
	\renewcommand{\footrulewidth}{0.0 pt}
% L:left, R:right, C:center, O:odd, E:even.
\fancyhf{} % Debe estar aquí para funcionar la configuración
% Encabezado superior:
\fancyhead[RO,LE]{}
\fancyhead[RE,LO]{}
% Encabezado inferior:
\fancyfoot[L]{}
\fancyfoot[C]{\thepage} % Número de páginas
\fancyfoot[R]{}
}

\newcommand{\EncYpiepaGcuerpo}{
% Estilo:
% ***	Definición de la página inicial de cada capítulo:
% Configuración de la línea divisoria:
\fancypagestyle{plain}{
\renewcommand{\headrulewidth}{0.0 pt}
\renewcommand{\footrulewidth}{1.3 pt}
% Encabezado:
\fancyhead[L]{}
\fancyhead[C]{}
\fancyhead[R]{}
% Pié de página:
\fancyfoot[L]{Ingeniería en Computación}
\fancyfoot[C]{}
\fancyfoot[R]{Página \thepage \hspace{1 pt} de \pageref{LastPage}} % Número de páginas
}
% ***	Definición de las siguientes páginas:
\pagestyle{fancy}
% Configuración de la línea divisoria:
	\renewcommand{\headrulewidth}{1.3 pt}
	\renewcommand{\footrulewidth}{1.3 pt}
% L:left, R:right, C:center, O:odd, E:even.
\fancyhf{} % Debe estar aquí para funcionar la configuración
% Encabezado:
\fancyhead[LE]{\slshape\nouppercase{\leftmark}} % \chapter{•}
\fancyhead[RO]{\slshape\nouppercase{\rightmark}} % \section{•}
%
\fancyhead[C]{}
\fancyfoot[L]{}
\fancyfoot[C]{}
\fancyfoot[RO,LE]{\thepage}
}

%**************************************************************************************************

\newcommand{\EncYpiepaGepilogo}[1]{
\pagestyle{fancy}
% Configuración de la línea divisoria:
	\renewcommand{\headrulewidth}{1.3 pt}
	\renewcommand{\footrulewidth}{1.3 pt}
% L:left, R:right, C:center, O:odd, E:even.
\fancyhf{} % Debe estar aquí para funcionar la configuración
% Encabezado:
\fancyhead[R]{}
\fancyhead[L]{}
\fancyhead[LE]{\slshape\nouppercase{#1}} % \chapter{•}
}

%%%%%%%%%%%%%%%%%%%%%%%%%%%%%%%%%%
%%%% PREÁMBULO PARA ACRÓNIMOS %%%%%
%%%%%%%%%%%%%%%%%%%%%%%%%%%%%%%%%%%%

\usepackage[acronym,shortcuts,nonumberlist,toc,nomain]{glossaries}

% acronym --> se diferencia de los glosarios, que son "definiciones"
% nonumberlist --> no genera numeración
% toc --> aparece en la tabla de contenidos
% nomain --> no genera el archivo .glo, sino que toma los acrónimos del .tex generado manualmente
\setglossarystyle{long}					% Utiliza el estilo "long" (el "super" pisa el pié de página)
\renewcommand{\glsgroupskip}{}			% Hace equidistantes los espacios entre acronimos
\renewcommand{\glsnamefont}[1]{\textbf{#1}\vspace{.3cm}}	% Pone en negrita los acrónimos
	\setlength{\glsdescwidth}{4.67 in}	% Ancho que ocupa la Descripción (8.27"-1.6"-1"-1"deACR)
\makeglossaries							% Invoca codificación de carácteres, soporte de idiomas y xindy ó makeindex
\usepackage{tabularx}

%%%%%%%%%%%%%%%%%%%%%%%%%%%%%%%%%%%%%%%%%%%
%%%% PREÁMBULO PARA DIAGRÁMA SINÓPTICO %%%%%
%%%%%%%%%%%%%%%%%%%%%%%%%%%%%%%%%%%%%%%%%%%%%

\newenvironment{subgroup}
  {$\left\{\tabular{l}}
  {\endtabular\right.$}

%%%%%%%%%%%%%%%%%%%%%%%%%%%%%%%%%%%%%%%%%%%
%%%% PAQUETE PARA COMENTARIOS %%%%%
%%%%%%%%%%%%%%%%%%%%%%%%%%%%%%%%%%%%%%%%%%%
\usepackage{verbatim} % comentarios
%%%%%%%%%%%%%%%%%%%%%%%%%%%%%%%%%%%%%%%%%%%
%%%% AGREGAR UNA PAGINA EN BLANCO %%%%%%
\usepackage{afterpage}
%%%%%%%%%%%%%%%%%%%%%%%%%%%%%%%%%%%%%%%%%%%